\section{讨论}

本文的主要价值，不在于再增加一组新的实验结果，而在于把已经完成的大量结果重新组织为一条适合临床期刊的证据链。按照“原始临床基线 -> 传统影像组学 -> 多模态深度融合 -> 多架构与消融 -> 多源融合扩展验证”的顺序重构后，可以得到三个更清晰的结论。

第一，\textbf{原始临床信息仍然是当前任务中最稳定的基础信号。}直接从原始术前变量构建的 Logistic 回归在 CentralLNM 上达到 AUC 0.6859$\pm$0.0658，提示临床变量本身具有较强可复核性和解释性，应作为主文参照基线长期保留。

第二，\textbf{影像信息确实有价值，但其临床意义更体现在“增量价值”而非“替代价值”。}传统 ROI Radiomics-only 表现偏弱，说明手工统计特征难以单独承担主要判别任务；Clinical+ROI Radiomics 虽较原始临床略有提高，但增益仅为 0.0064，且折间波动未见改善。真正带来相对稳定收益的，是患者级 MIL、临床主导门控与先验约束共同组成的深度融合框架：其相对于原始临床的 AUC 增益为 0.0245，同时标准差由 0.0658 降至 0.0253。换言之，本研究支持的并不是“深度学习一定优于临床”，而是“当影像分支被组织成更符合临床逻辑的结构时，才能以有限但更稳定的方式成为临床的补充”。

第三，\textbf{多源融合当前更像潜在方向，而不是已经完成定型的最终方案。}在 ViT-B/16 协议下，引入 ROI radiomics 的 V2 版本并未超过 V1（0.6404 vs 0.7104），且已回落至原始临床基线以下，说明“信息源增加”并不必然带来增益。该结果提示后续应优先优化组学分支与深度表征之间的交互机制，而非继续扩大输入维度。

本研究也有若干局限。其一，全部结果均来源于单中心回顾性数据与内部 5-fold 验证，尚缺乏独立外部验证。其二，虽然本文已尽量统一公平比较协议，但不同部分的历史实验在骨干网络、融合器和辅助指标上仍有一定差异，因此不宜将所有表格做机械的逐项横向比较。其三，主文为保证期刊可读性，刻意压缩了大量过程性排查和敏感性分析；这有助于形成清晰主线，但也意味着某些工程细节需要读者结合补充材料理解。

从临床角度看，本文最稳妥的结论是：CentralLNM 的术前评估不应脱离临床变量，而更适合采用“临床主导、影像补充”的融合策略。对于当前数据而言，完整深度融合模型已经具备作为后续外部验证候选模型的条件；而下一阶段真正值得提升的，不是继续堆叠更多模态或骨干，而是围绕层级结构建模、患者级弱监督可解释 MIL 以及个体化证据融合，证明影像在哪些患者上能够改变原始临床决策。
